%!TEX root = ../thesis.tex
\chapter{Tissue Preprocessing Methodology Checks} \label{app:tissue_methodology}

Following the locked tissue pipeline in Section~\ref{sec:preprocessing} and Section~\ref{sec:patch_extraction}, \textit{FINAL Step~2.7 --- Tissue Preprocessing Methodology Checks} runs five sensitivity analyses on the 434 background patches under the same image-level five-seed WS protocol as the main tissue replication (Subspace $k$-NN ensemble; seeds $\{1,7,34,42,99\}$). Outputs are saved under \texttt{data/outputs/tissue\_preprocessing\_methodology\_checks/}. These checks do \emph{not} change any headline result in Chapters~\ref{sec:results}--\ref{sec:discussion}; they document that alternative preprocessing and sampling choices do not outperform the adopted pipeline.

\begin{table}[h]
\centering
\caption[Tissue preprocessing sensitivity checks]{Tissue preprocessing sensitivity checks (FINAL Step~2.7). WS AUROC values are five-seed means on the locked evaluation protocol.}
\label{tab:tissue_methodology_checks}
\small
\begin{tabular}{p{3.2cm}p{4.8cm}p{4.6cm}}
\toprule
\textbf{Check} & \textbf{Procedure} & \textbf{Outcome} \\
\midrule
Raw-coordinate pixel sources &
Re-crop patches at Step~1.3 coordinates from raw DICOM (min-max or percentile scaling, with/without CLAHE) &
Current preprocessed patches remain best (AU-ROC $0.895$); raw min-max is clearly worse ($0.812$); no-CLAHE / raw percentile match at $0.881$ \\
Otsu-selection audit &
Compare accepted patches to approximate rejected candidates (same images; fail the $\ge$65\% target-tissue rule; $n=60$ per class) &
Largest shifts are in \texttt{target\_frac} and patch size, as expected from the purity criterion; smaller shifts in location and high-band FFT \\
No-CLAHE control &
WS on percentile-scaled, masked patches without CLAHE &
AU-ROC $0.881$ vs.\ $0.895$ for the locked pipeline (${\sim}$1.4~pp loss) \\
Unsupervised WS sanity &
$k=2$ KMeans on scaled WS features (no classifier) &
Adjusted Rand ${\sim}0.22$--$0.29$ vs.\ ACR labels; supervised WS performance is unchanged \\
Frequency diagnostic &
Radial FFT power: fatty vs.\ fibroglandular per pixel source &
Locked pipeline shows the largest mean spectral gap ($0.459$ vs.\ $0.27$--$0.38$ for alternatives) \\
\bottomrule
\end{tabular}
\end{table}

\textbf{Interpretation.} The percentile-clip, CLAHE, and Otsu/ACR-guided patch-sampling choices adopted in the main pipeline are internally consistent: they match or beat raw-coordinate alternatives on WS AUROC and preserve more fatty/fibroglandular spectral separation than min-max or no-CLAHE paths. The Otsu audit confirms that accepted patches differ systematically from failed candidates---chiefly in target-tissue purity and window size---because the sampler enforces an explicit purity rule; this is intentional task definition, not evidence of test-set leakage (evaluation remains image-level held-out). The unsupervised clustering check shows that WS features do not form two clean clusters matching ACR labels without supervision; this does not contradict the supervised WS baseline. No CNN or fusion branch was retrained on alternate pixel sources, since WS did not identify a better candidate worth the additional compute.
